\section{Strong-Bisimulation Quotients: Well-Founded and Labelled Cases}
\label{sec:bisimulation}

A quotient may contain a cycle even when its source graph does not: different
quotient edges can be witnessed by incompatible representatives.  Successor
matching prevents that defect whenever the source admits no infinite forward
path.

\begin{theorem}[Forward-well-founded bisimulation obstruction]
\label{thm:well-founded-bisim}
Let $G=(V,E)$ be forward well-founded and let $\sim$ be a strong forward
bisimulation.  Then the quotient graph $G/{\sim}$ is acyclic.
\end{theorem}

\begin{proof}
Assume, for contradiction, that the quotient contains a directed cycle
\[
  C_0\longrightarrow C_1\longrightarrow\cdots
  \longrightarrow C_{m-1}\longrightarrow C_m=C_0,
  \qquad m\geq1.
\]
For each $i\in\{0,\ldots,m-1\}$, choose representatives
$a_i\in C_i$ and $b_{i+1}\in C_{i+1}$ such that $a_iEb_{i+1}$, with indices
read modulo $m$.

Choose any $v_0\in C_0$.  Suppose $v_j\in C_i$ has been chosen, where
$i\equiv j\pmod m$.  Since $v_j\sim a_i$ and $a_iEb_{i+1}$, successor
matching supplies $v_{j+1}$ such that
\[
  v_jEv_{j+1}
  \qquad\text{and}\qquad
  v_{j+1}\sim b_{i+1}.
\]
Thus $v_{j+1}\in C_{i+1}$.  Induction constructs an infinite directed path
$v_0Ev_1Ev_2E\cdots$ in $G$, contradicting forward well-foundedness.
\end{proof}

\begin{corollary}[Finite-DAG case]
\label{cor:finite-bisim}
A strong forward-bisimulation quotient of a finite directed acyclic graph is
acyclic.
\end{corollary}

\begin{proof}
An infinite path in a finite graph repeats a vertex, and the positive segment
between two repetitions is a directed cycle.  A finite DAG is therefore
forward well-founded, so \cref{thm:well-founded-bisim} applies.
\end{proof}

Acyclicity alone cannot replace forward well-foundedness for an arbitrary
infinite source.

\begin{remark}[Infinite-ray counterexample]
\label{rem:infinite-ray}
Let $V=\N$ and put one edge $n\to n+1$ at every vertex.  This infinite graph
is acyclic but not forward well-founded.  The universal equivalence relation,
under which all vertices are equivalent, is a strong forward bisimulation:
the unique successors of any two vertices are again equivalent.  Its quotient
has one vertex and a self-loop.  Thus \cref{thm:well-founded-bisim} cannot be
enlarged to all infinite DAGs by retaining acyclicity alone.
\end{remark}

For the wheel system, an independent grading argument applies when one
state-class decoder must be exact for every representative of that class.

\begin{proposition}[Level-injective labels preserve acyclicity]
\label{prop:label-grading}
Let $V=\bigsqcup_{k\geq0}V_k$, and suppose every edge goes from $V_k$ to
$V_{k+1}$.  Let $\lambda:V\to\Lambda$ be constant on each $V_k$, with
pairwise distinct values across levels.  If
\begin{equation}
  v\sim w \quad\Longrightarrow\quad \lambda(v)=\lambda(w),
  \label{eq:label-respecting}
\end{equation}
then $G/{\sim}$ is acyclic.  Graph finiteness and bisimulation are not needed.
\end{proposition}

\begin{proof}
Condition \eqref{eq:label-respecting} and injectivity of the level labels
force every equivalence class to lie in a single $V_k$.  Define
\[
  \overline{\level}([v])=k \qquad\text{when }v\in V_k.
\]
This quotient grading is well-defined.  If $[v]\to[w]$, some representatives
$v'\in[v]$ and $w'\in[w]$ satisfy $v'Ew'$, hence
\[
  \overline{\level}([w])=\overline{\level}([v])+1.
\]
The quotient level increases along every positive-length path, so no path can
return to its starting class.
\end{proof}

\begin{corollary}[Representative-exact next-prime decoder]
\label{cor:q-label}
Let $G$ be the graded wheel graph and suppose there is one state-class decoder
$D:V/{\sim}\to\N$ satisfying
\begin{equation}
  D([v])=q_{\level(v)+1}
  \qquad\text{for every source representative }v.
  \label{eq:class-decoder-exact}
\end{equation}
Then the quotient cannot merge different levels and cannot contain a directed
cycle.
\end{corollary}

\begin{proof}
If $v\sim w$, then \eqref{eq:class-decoder-exact} gives
$q_{\level(v)+1}=D([v])=D([w])=q_{\level(w)+1}$.  By
\cref{lem:prime-enumeration}, the values $q_k$ are pairwise distinct, so
$\level(v)=\level(w)$.  Apply \cref{prop:label-grading} with
$\lambda(v)=q_{\level(v)+1}$.
\end{proof}

\begin{remark}[What the quotient results do not cover]
The proof of \cref{thm:well-founded-bisim} uses successor matching and the
absence of an infinite forward source path.  It does not apply to an infinite
ray, a one-way simulation, a bounded-radius observational equivalence, or an
arbitrary graph homomorphic image.  The proof of
\cref{prop:label-grading} applies only when the state-class decoder is exact
for every representative.  An edge decoder or a path-window decoder that
distinguishes representatives within one state class falls outside that
proposition.  Falling outside either theorem creates an open proof
obligation, not evidence of an arithmetic periodic orbit.
\end{remark}

For a finite cutoff, \cref{cor:finite-bisim} removes the proposed positive
branch entirely: a correctly computed strong-bisimulation quotient cannot
produce a cycle.  Running the quotient code could still serve as a regression
test, but it could not supply positive Stage-02 evidence.
